% ============================================================
%  PRODUCT BREAKDOWN STRUCTURE (PBS) — Phase 1 : Démarrage
%  Time'Eats — Cellule PMO
% ============================================================
\documentclass[11pt,a4paper,oneside]{report}
\usepackage{../style/consulting}
\hypersetup{pdftitle={Product Breakdown Structure (PBS) — Time'Eats PMO}}
\pagestyle{fancy}\fancyhf{}
\renewcommand{\headrulewidth}{0pt}
\fancyhead[L]{\begin{tikzpicture}[remember picture,overlay]
  \draw[cnNavy,line width=0.5pt](0,0)--(\textwidth,0);
\end{tikzpicture}\vspace{2pt}\timeeatslogo\hspace{4pt}\small\textcolor{cnSlate}{Time'Eats \textbf{·} Product Breakdown Structure (PBS)}}
\fancyhead[R]{\small\textcolor{cnSlate}{\thepage}}
\fancyfoot[C]{\tiny\textcolor{cnSlate}{Document confidentiel --- Cellule PMO --- CESI MAALSI 2026}}
\begin{document}\small

\begin{center}
{\LARGE\bfseries\color{cnNavy} Product Breakdown Structure (PBS)}\\[4pt]
\textcolor{cnOrange}{\rule{10cm}{1pt}}\\[4pt]
{\large\color{cnSlate} Phase 1 --- Démarrage}
\end{center}

\vspace{6pt}
\renewcommand{\arraystretch}{1.4}
\begin{tabularx}{\textwidth}{L{3.5cm} Y L{3.5cm} Y}
\toprule
\rowcolor{cnNavy}\th{Projet} & \th{} & \th{Date} & \th{} \\
\midrule
Nom        & \textit{[Nom du projet]} & Rédacteur & \textit{[Prénom NOM]} \\
\rowcolor{cnIce}
Code       & \textit{[CODE-XX]}       & Version   & 1.0 \\
\bottomrule
\end{tabularx}

\section*{1 --- Objet du document}

Le \textbf{Product Breakdown Structure (PBS)} décompose hiérarchiquement le \textbf{produit livré} en sous-systèmes, composants et éléments. Il complète le WBS (P2-06) qui décompose le \textit{travail} à réaliser :

\begin{itemize}
  \item \textbf{PBS} = \textit{quoi est livré} (vue produit, statique, orientée client / utilisateur).
  \item \textbf{WBS} = \textit{comment le travail est organisé} (vue effort, dynamique, orientée équipe projet).
\end{itemize}

Le PBS est produit \textbf{avant} le WBS : il sert de base à la rédaction du Cahier des Charges (P1-02ter), à l'estimation budgétaire (P2-08), au découpage en lots de travail (P2-06) et à la définition des critères de recette (P5-19 / P5-20).

\section*{2 --- Décomposition produit (4 niveaux)}

\noindent
\textbf{Niveau 0 :} \textit{[Nom du produit final]} \\
\textbf{Niveau 1 :} sous-systèmes (briques fonctionnelles majeures) \\
\textbf{Niveau 2 :} composants (éléments livrables identifiables) \\
\textbf{Niveau 3 :} éléments unitaires (modules, services, écrans, livrables documentaires)

\vspace{6pt}
\renewcommand{\arraystretch}{1.3}
\begin{tabularx}{\textwidth}{C{1.4cm} L{4cm} L{4cm} Y}
\toprule
\rowcolor{cnNavy}\th{Code} & \th{Niveau 1 (sous-système)} & \th{Niveau 2 (composant)} & \th{Niveau 3 (éléments)} \\
\midrule
1   & \textit{[Sous-système 1]} & & \\
\rowcolor{cnIce}
1.1 &                            & \textit{[Composant 1.1]} & \textit{[Éléments]} \\
1.2 &                            & \textit{[Composant 1.2]} & \textit{[Éléments]} \\
\midrule
2   & \textit{[Sous-système 2]} & & \\
\rowcolor{cnIce}
2.1 &                            & \textit{[Composant 2.1]} & \textit{[Éléments]} \\
2.2 &                            & \textit{[Composant 2.2]} & \textit{[Éléments]} \\
\midrule
3   & \textit{[Sous-système 3]} & & \\
\rowcolor{cnIce}
3.1 &                            & \textit{[Composant 3.1]} & \textit{[Éléments]} \\
\bottomrule
\end{tabularx}

\section*{3 --- Vue arborescente synthétique}

\begin{verbatim}
[Produit final]
+-- 1. [Sous-systeme 1]
|   +-- 1.1 [Composant]
|   +-- 1.2 [Composant]
+-- 2. [Sous-systeme 2]
|   +-- 2.1 [Composant]
|   +-- 2.2 [Composant]
+-- 3. [Sous-systeme 3]
    +-- 3.1 [Composant]
\end{verbatim}

\section*{4 --- Correspondance PBS $\leftrightarrow$ autres documents}

\renewcommand{\arraystretch}{1.3}
\begin{tabularx}{\textwidth}{L{4cm} Y}
\toprule
\rowcolor{cnNavy}\th{Document amont / aval} & \th{Lien avec le PBS} \\
\midrule
P1-02 Note de cadrage      & Le PBS détaille le périmètre \textit{Inclus} formulé dans la note de cadrage. \\
\rowcolor{cnIce}
P1-02ter CDC               & Chaque composant du PBS donne lieu à une ou plusieurs exigences fonctionnelles. \\
P2-06 WBS                  & Chaque composant du PBS est associé à un ou plusieurs lots de travail (WP) du WBS. \\
\rowcolor{cnIce}
P2-08 Budget               & Le PBS sert de grille de répartition des coûts par sous-système. \\
P5-19 PV recette fonct.    & Les critères d'acceptation sont rattachés aux éléments de niveau 3 du PBS. \\
\bottomrule
\end{tabularx}

\begin{reco}
\textbf{Bonnes pratiques :} (i) limiter le PBS à 4 niveaux pour rester lisible ; (ii) ne pas confondre PBS et WBS --- un PBS qui devient verbal (\textit{« réaliser », « tester », « livrer »}) doit être réécrit en termes de \textit{produit} ; (iii) faire valider le PBS par la MOA \textbf{avant} la rédaction du CDC.
\end{reco}

\end{document}
